\noindent\textbf{Resultado.} $h(t)=24-\dfrac4{15}t$ cm para $0\le t<90$ días; el vaciado se alcanza en $t=90$ días.

\noindent\textit{Desarrollo.} La sección horizontal de la esfera a altura h tiene radio cuadrado $16^2-(h-16)^2=h(32-h)$, de modo que su área es $A(h)=\pi h(32-h)$. También $dV/dh=32\pi h-\pi h^2=A(h)$, como debe ocurrir al añadir una lámina de espesor dh. Si k>0 es la constante de evaporación, con unidades cm/día, el balance es $dV/dt=-kA(h)$. Por la regla de la cadena, $A(h)h^{\prime}=-kA(h)$. Mientras $0<h<32$, A es positiva y $h^{\prime}=-k$. Integrar da $h=24-kt$. El dato $h(15)=20$ fija $k=4/15$ cm/día. La solución conserva ambas alturas y el balance de volumen; llega a h=0 en 90 días. En ese extremo A=0 y la cancelación anterior ya no procede; prolongar la recta produciría alturas negativas ajenas al modelo.
